% Literature Review, section 3, with the summary table. Notes for the whole chapter are in
% 00_chapter_intro.tex.
\section{Geometry and Risk of Coronary Disease}
\label{sec:lit-risk}

Estimating an individual's risk of developing plaque from geometry requires geometry measured before
plaque appears and a later scan showing whether it does.
Geometry measured once plaque is present cannot show which came first, as an artery remodels while
plaque grows \cite{glagov1987remodeling}.

Studies with follow-up have mostly concerned disease that was already present.
Sakellarios et al.\ (2017) \cite{sakellarios2017bifurcation} followed bifurcations with repeated
CCTA in a small group of patients with established disease, reporting that WSS predicted lumen
reduction best when the daughter branch was included.
Han et al.\ (2022) \cite{han2022plaque} compared a small sample of patients who later had an acute
coronary syndrome with matched controls from a registry of patients referred for suspected disease.
Their findings indicate that the lesions that caused the event lay nearer the ostium and more often
at bifurcations than other lesions in the same patients.
However, the typical event followed within weeks of the scan, so these lesions were already present
at imaging.
Tommasino et al.\ (2024) \cite{tommasino2024clap} followed patients referred for CCTA for one year,
linking a wide left main bifurcation angle to progression of proximal LAD disease in a cohort in which
many already had proximal LAD disease.
Most of these events were progression on imaging requested by clinicians, and their reported
intervals cannot be reproduced.
In contrast, Won et al.\ (2022) \cite{won2022paradigm} followed registry patients without plaque on
a first CCTA to a repeat scan some years later, observing new plaque in a substantial share of them
and relating plaque progression to age, sex, hypertension, diabetes, hyperlipidemia, obesity,
smoking and HbA1c but not to geometry.

Despite these advancements, a significant gap remains in the current literature.
Cohorts suited to this question have not measured geometry, since the plaque-free registry of Won
et al.\ included none and the random population sample of Bergstr\"om et al.\
\cite{bergstrom2021scapis} was scanned once and so had no follow-up.
No study identified in this search has combined automated measurement of the whole coronary tree with
follow-up to new plaque in a population selected at random rather than referred for suspected
disease.
An overview of the studies relating coronary geometry to disease is given in
Table~\ref{tab:literature}.

\begin{table}[h]
  \centering
  \footnotesize
  \begin{tabular}{@{}>{\raggedright\arraybackslash}p{0.16\linewidth}%
                  >{\raggedright\arraybackslash}p{0.22\linewidth}%
                  >{\raggedright\arraybackslash}p{0.25\linewidth}%
                  >{\raggedright\arraybackslash}p{0.29\linewidth}@{}}
    \toprule
    Author \& Year & Cohort & Geometric Approach & Main Findings \\
    \midrule
    Zhang (2024) & 39 left coronary trees, stenosed and stenosis-free
      & Whole-tree curvature on CCTA, cross-sectional
      & Curvature separates stenosed from stenosis-free trees. \\
    Li (2011) & 1010 consecutive angiography patients
      & Bend count on invasive angiograms, cross-sectional
      & Tortuosity inversely associated with CAD. \\
    Tello Ayala (2026) & 22,334 referred patients
      & Turning-angle profile along the RCA, cross-sectional
      & Tortuosity higher with CAD; the full profile identifies CAD better than its mean. \\
    Han (2022) & 116 referred patients with later ACS, 116 controls
      & Lesion position and shape on CCTA, follow-up to ACS
      & Event-causing lesions nearer the ostium and at bifurcations; events within weeks. \\
    \bottomrule
  \end{tabular}
  \caption{Overview of studies relating coronary geometry to disease, highlighting the cohort, the
    geometric approach and the main findings. None measured geometry in a randomly selected
    population and followed it to new plaque. ACS: acute coronary syndrome; CAD: coronary artery
    disease; CCTA: coronary computed tomography angiography; RCA: right coronary artery.}
  \label{tab:literature}
\end{table}
